% !TeX root = ../main.tex

% Edit this title for each week's entry.
\section{Week 10 - Inflation}

\topic{Introduction, Causes and Effects}

\begin{example}[Reflection: Hamburger Prices]
    How much does a McDonald's hamburger cost?
    \\\db{On an early McDonald's menu, a hamburger was \dbb{15\textcent}. The Canadian Big Mac rose from about \$1.90 in 1986 to over \$6.50 by 2018. Average Canadian undergraduate tuition rose from about \$4300 to \$6400 between 2006 and 2017. The same goods cost more dollars over time.}
\end{example}

\begin{definition}[Inflation]
    A rise in the \dbb{average prices} of goods and services over time. It
    reflects the decline in the \dbb{purchasing power} of the dollar over
    time. \rb{Deflation} is the opposite.
\end{definition}

\heading{What Causes Inflation?}
\begin{itemize}
    \item \rb{Demand pull}: demand rises faster than supply
    \item \rb{Cost push}: production costs (wages, materials, energy) rise and are passed on
    \item \rb{Expanding money supply}: more money chasing the same goods
    \item \rb{Expectations}: if people expect prices to rise, they demand higher wages and raise prices now
\end{itemize}

\heading{What Are the Effects of Inflation?}
\begin{itemize}
    \item \dbb{Benefits borrowers}: they repay with dollars worth less
    \item \dbb{Hurts} lenders, savers, the unemployed, retirees (fixed incomes), etc.
    \item Menu costs (changing prices) and tax distortions
    \item In the \b{short run only}, higher inflation can lead to lower unemployment
    \item At low levels it can stimulate the economy; at high levels (\rb{hyperinflation}) it can be a disaster
\end{itemize}

\heading{Bank of Canada Inflation Target}
The primary mission of the Bank of Canada is to control inflation, with a
target of \dbb{2\%} per year.

\begin{example}[Reflection: Why Not 0\%?]
    Why don't we target 0\% inflation?
    \\\db{A target of 0\% leaves no buffer against \dbb{deflation}. Falling prices make people delay spending and make debts heavier in real terms, which can spiral into a recession. A small positive rate also gives the central bank room to cut interest rates, and lets real wages adjust without cutting anyone's nominal pay.}
\end{example}

\begin{remark}
    Remember that the interest rate is interest / principal!
\end{remark}

\heading{Summary}
\begin{enumerate}
    \item Inflation $=$ rising prices
    \item Caused by rising demand, increasing costs, an expanding money supply, or expectations
    \item OK at low levels; a disaster at high levels!
\end{enumerate}

\topic{Measuring Inflation: Consumer Price Index}

\begin{definition}[Consumer Price Index (CPI)]
    The most commonly used inflation index: the \dbb{weighted average of
    prices} of a basket of goods and services of primary consumer needs
    (food, housing, transportation, apparel, medical care, recreation,
    education, \dots). It relates the average price of a \b{fixed}
    \rb{basket} of goods in a given period to its price in a \rb{base
    period}:
    \[\boxed{\text{CPI}_t = \frac{\text{Cost of basket in year } t}{\text{Cost of basket in base year}} \times 100}\]
    The current CPI base year is 2002, with index 100. The index for any
    other year is the number of dollars needed that year to buy the basket
    that cost \$100 in 2002.
\end{definition}

\begin{center}
\begin{tabular}{cc@{\hspace{3em}}cc}
    \toprule
    Year & CPI & Year & CPI \\
    \midrule
    2002 & 100.0 & 2011 & 119.9 \\
    2003 & 102.8 & 2012 & 121.7 \\
    2004 & 104.7 & 2013 & 122.8 \\
    2005 & 107.0 & 2014 & 125.2 \\
    2006 & 109.1 & 2015 & 126.6 \\
    2007 & 111.5 & 2016 & 128.4 \\
    2008 & 114.1 & 2017 & 130.4 \\
    2009 & 114.4 & 2018 & 133.4 \\
    2010 & 116.5 & 2019 & 136.0 \\
    \bottomrule
\end{tabular}
\\[2pt] \textit{\small Canadian historical CPI (Statistics Canada)}
\end{center}

\begin{example}[Calculating CPI]
    \begin{center}
    \begin{tabular}{llrrrr}
        \toprule
        & & Hamburger & Jeans & Movie ticket & Total \\
        \midrule
        2002 & Price & \$4.00 & \$70.00 & \$6.00 & \\
        & Quantity & 40 & 1 & 5 & \\
        & Expenditure & \db{\$160} & \db{\$70} & \db{\$30} & \db{\$260} \\
        \midrule
        2020 & Price & \$6.00 & \$80.00 & \$8.00 & \\
        & Quantity & 40 & 1 & 5 & \\
        & Expenditure & \db{\$240} & \db{\$80} & \db{\$40} & \db{\$360} \\
        \bottomrule
    \end{tabular}
    \end{center}
    \[\text{CPI}_{2002} = \frac{260}{260} \times 100 = 100, \qquad \text{CPI}_{2020} = \frac{360}{260} \times 100 = \db{138.5}\]
    \[\text{Inflation 2002--2020} = \frac{138.5 - 100}{100} = \db{38.5\%}\]
\end{example}

\heading{Calculating Inflation from CPI}
\[\boxed{1 + f_t = \frac{\text{CPI}_t}{\text{CPI}_{t-1}}, \qquad (1 + f_a)^n = \frac{\text{CPI}_{t+n}}{\text{CPI}_t}}\]
where $f_t$ is the inflation rate in year $t$ and $f_a$ is the average annual
inflation rate over $n$ years.

\begin{example}[Inflation in 2019]
    \[f_{2019} = \frac{\text{CPI}_{2019}}{\text{CPI}_{2018}} - 1 = \frac{136.0 - 133.4}{133.4} = \db{1.95\%}\]
\end{example}

\begin{example}[Reflection: Average Annual Inflation]
    The CPI was 102.8 in 2003 and 126.6 in 2015. What is the average annual
    inflation rate from 2003 to 2015?
    \[f_{2003\text{--}2015} = \frac{126.6}{102.8} - 1 = 23.2\%\]
    \db{This is spread over $2015 - 2003 = 12$ years of compounding:}
    \[(1 + f_a)^{12} = 1.232 \;\Rightarrow\; f_a = \db{1.75\%}\]
    \db{(The handwritten slide uses an exponent of $2015 - 2002 = 13$ and gets 1.6\%. From 2003 to 2015 is only 12 years.)}
\end{example}

\heading{Summary}
\begin{enumerate}
    \item The CPI is built from the price changes of a fixed basket of goods
    \item Inflation is determined by changes in the CPI
\end{enumerate}

\topic{Calculations with Inflation}

\begin{example}[Investing with Inflation]
    A \$100 investment pays 5\% interest for 3 years, with inflation
    expected to be 2\% per year. What is the value of the investment? What
    is its \emph{real value} after 3 years?
    \begin{center}
    \begin{tabular}{lcccc}
        \toprule
        Year & 0 & 1 & 2 & 3 \\
        \midrule
        Investment value & \$100 & \$105 & \$110.25 & \$115.76 \\
        Price of good & \$100 & \$102 & \$104.04 & \$106.12 \\
        Purchasing power ratio & 1.00 & 1.03 & 1.06 & 1.09 \\
        Real value (today's \$) & \$100 & \$102.94 & \$105.97 & \db{\$109.09} \\
        \bottomrule
    \end{tabular}
    \end{center}
    \[100(1.05)^3 = \db{\$115.76}, \qquad \text{real value} = 100\left(\frac{1.05}{1.02}\right)^3 = \db{\$109.09}\]
    \db{The money grows at 5\%, but its purchasing power only grows at $1.05/1.02 - 1 = 2.94\%$ per year, the \dbb{real} rate.}
\end{example}

\begin{definition}[Actual and Real Dollars]
    \begin{itemize}
        \item \rb{Actual (current, nominal) dollars}: expressed in the monetary units at the time the cash-flows occur
        \item \rb{Real (constant) dollars}: expressed in monetary units of \dbb{constant purchasing power}. They must always be associated with a particular date (e.g.\ ``2026 dollars'')
    \end{itemize}
\end{definition}

\begin{definition}[Actual and Real Interest Rates]
    \begin{itemize}
        \item \rb{Actual interest rate} ($i$ or $i_A$): the observed interest rate, based on actual dollars
        \item \rb{Real interest rate} ($i'$ or $i_R$): the interest rate based on dollars of constant purchasing power
    \end{itemize}
    \[\boxed{\begin{gathered}
        1 + i_R = \frac{1 + i_A}{1 + f}\\[4pt]
        \begin{array}{@{}l@{}}
            i_A: \text{actual interest rate}\\
            i_R: \text{real interest rate}\\
            f: \text{inflation rate}
        \end{array}
    \end{gathered}}\]
    \db{For small rates, $i_R \approx i_A - f$.}
\end{definition}

\heading{Economic Valuation with Inflation}
\begin{center}
\begin{tabular}{lll}
    \toprule
    Cash-flows in & Adjust for inflation? & Discount with \\
    \midrule
    \rb{Actual} dollars & Yes (they include inflation) & \dbb{Actual} MARR \\
    \rb{Real} dollars & No & \dbb{Real} MARR \\
    \bottomrule
\end{tabular}
\end{center}

\begin{remark}
    Most market interest rates are \textbf{actual} rates: loan interest
    rates, bond yields, expected stock returns (CAPM), and contracts (unless
    explicitly stated otherwise, e.g.\ defined benefit pension plans).
\end{remark}

\db{Historically (Canada, 1958--2009), long-term government bond yields have tracked the inflation rate, staying above it except during the mid-1970s inflation peak.}

\begin{example}[Cash-flow with Inflation]
    A company invests \$100k in a new machine that brings benefits of \$30k
    per year for 5 years. It expects a \b{real} return of 6\%, and inflation
    is 3\%.
    \begin{enumerate}[label = \alph*)]
        \item If the dollar values are \b{actual}, find the PW.
        \\\db{Actual dollars need the actual MARR:}
        \[i_A = (1.06)(1.03) - 1 = 9.18\%\]
        \[\text{PW} = -100\text{k} + 30\text{k}(P/A,9.18\%,5) = \db{\$16.1\text{k}}\]
        \item If the dollar values are \b{real}, find the PW.
        \[\text{PW} = -100\text{k} + 30\text{k}(P/A,6\%,5) = \db{\$26.4\text{k}}\]
        \db{Real benefits keep their purchasing power (they grow with inflation in actual dollars), so the project is worth more.}
    \end{enumerate}
\end{example}

\begin{example}[Loan with Inflation]
    A bank quotes 5\% compounded semi-annually. If you expect inflation of
    2\%, what is the real effective rate of interest?
    \[1 + i_{A,\text{eff}} = \left(1 + \frac{0.05}{2}\right)^2 \;\Rightarrow\; i_{A,\text{eff}} = 5.0625\%\]
    \[i_{R,\text{eff}} = \frac{1.050625}{1.02} - 1 = \db{3.002\%}\]
\end{example}

\begin{example}[Bond with Inflation]
    A 10-year bond with 5 years left to maturity has a 4\% coupon, paid
    semi-annually. Inflation is expected to average 3\% over the next 5
    years. If you target a real (effective) return of 5\%, how much should
    you pay for this bond (per \$100 face)?
    \\\db{Bond cash-flows are fixed in actual dollars, so convert to an actual rate, then to a semi-annual rate:}
    \[i_A = (1.05)(1.03) - 1 = 8.15\%, \qquad i_s = (1.0815)^{1/2} - 1 = 3.995\%\]
    \[P = 2(P/A,3.995\%,10) + 100(P/F,3.995\%,10) = \db{\$83.81}\]
\end{example}

\heading{Inflation with Tax Benefits}
In accounting we use actual values. Since CCA is calculated in \b{actual}
dollars, we \dbb{must} use the \b{actual} interest rate in the tax benefit
factors. It is convenient to write the factors as
\[\boxed{\text{CTF} = 1 - \tau_{1/2} = 1 - \frac{td(1 + i/2)}{(i + d)(1 + i)}, \qquad \text{CSF} = 1 - \tau_{db} = 1 - \frac{td}{i + d}}\]
\[\text{PW} = -FC \times \text{CTF} + S \times \text{CSF} \times (P/F,i,N)\]
where CTF is the \rb{capital tax factor} and CSF is the \rb{capital salvage
factor} (salvage treated as open books, $R = S$).

\begin{example}[Tax Benefit Factors with Inflation]
    A machine costs \$20,000, with an expected life of 7 years and a salvage
    value of \$6000 \b{in today's dollars}. The CCA rate is 20\% and taxes
    are 25\%. Find the PW, assuming a real MARR of 8\% and expected
    inflation of 4\%.
    \\\db{The tax benefit factors need the \dbb{actual} rate:}
    \[i_A = (1.08)(1.04) - 1 = 12.32\%\]
    \[\text{CTF} = 1 - \frac{0.25(0.2)(1.0616)}{(0.3232)(1.1232)} = 0.8538, \qquad \text{CSF} = 1 - \frac{0.25(0.2)}{0.3232} = 0.8453\]
    \db{Actual dollars: the salvage in year-7 dollars is $6000(1.04)^7$, discounted at the actual rate:}
    \[\text{PW} = -20{,}000(0.8538) + 6000\frac{(1.04)^7}{(1.1232)^7}(0.8453)\]
    \db{Real dollars: keep \$6000 and discount at the real rate, since $(1.04/1.1232)^7 = 1/(1.08)^7$:}
    \[\text{PW} = -20{,}000(0.8538) + \frac{6000}{(1.08)^7}(0.8453) = -17{,}076 + 2959 = \db{-\$14{,}116}\]
    \db{(The handwritten slide reports $-\$13{,}575$, but the terms it lists give $-\$14{,}116$.)}
\end{example}

\begin{remark}
    You \textbf{must} use the actual rate in the tax benefit factors, but
    you can use either actual or real dollars for the salvage value, as
    long as you discount with the matching rate.
\end{remark}

\begin{example}[Reflection: Machine with Declining Benefits]
    Find the PW of investing in a machine with annual (before-tax) benefits:
    \begin{center}
    \begin{tabular}{lr}
        \toprule
        Machine first cost & \$40,000 \\
        Expected life & 10 years \\
        Salvage value & \$5000 (real dollars) \\
        Annual benefits & \$10,000 in year 1, decreasing by 4\%/yr (real dollars) \\
        Actual MARR & 8\% \\
        CCA rate & 20\% \\
        Corporate tax rate & 30\% \\
        Expected inflation & 3\% \\
        \bottomrule
    \end{tabular}
    \end{center}
    \[\text{PW} = \text{PW}_{\text{machine}} + \text{PW}_{\text{benefits}}\]
    \db{\dbb{Machine} (tax factors at the actual rate):}
    \[\text{CTF} = 1 - \frac{0.3(0.2)(1.04)}{(0.28)(1.08)} = 0.7937, \qquad \text{CSF} = 1 - \frac{0.3(0.2)}{0.28} = 0.7857\]
    \[\text{PW}_{\text{machine}} = -40{,}000(0.7937) + 5000(0.7857)\frac{(1.03)^{10}}{(1.08)^{10}} = \db{-\$29{,}301}\]
    \db{\dbb{Benefits} are in real dollars with a real growth rate $g = -4\%$, so discount at the real rate:}
    \[i_R = \frac{1.08}{1.03} - 1 = 4.85\%\]
    \[(P/\text{Geom},r,g,n) = \frac{1 - \left(\frac{1 + g}{1 + r}\right)^n}{r - g}, \qquad \text{PW}_{\text{benefits}} = 10{,}000(P/\text{Geom},4.85\%,-4\%,10) = \db{\$66{,}199}\]
    \[\text{PW} = -29{,}301 + 66{,}199 = \db{\$36{,}898}\]
    \db{(The slide gets $\text{PW}_{\text{machine}} = -\$28{,}634$ and $\text{PW} = \$37{,}565$, but its listed terms give $-\$29{,}301$. Also, the slide does not tax the benefits. Since they are before-tax, they should be multiplied by $(1 - t)$: $\text{PW}_{\text{benefits}} = 0.7(66{,}199) = \$46{,}339$, which gives $\text{PW} = \dbb{\$17{,}038}$.)}
\end{example}

\begin{remark}
    If the growth rate were \textbf{actual} (and the \$10,000 still in real
    dollars), convert the first payment to actual dollars and discount
    everything at the actual rate:
    \[\text{PW}_{\text{benefits}} = 10{,}000(1.03)(P/\text{Geom},8\%,-4\%,10)\]
\end{remark}

\heading{Summary}
\begin{enumerate}
    \item \b{Actual dollars}: monetary units at the time the cash-flows occur. \b{Real dollars}: constant purchasing power, tied to a particular date
    \item The \b{actual} interest rate is the observed rate; the \b{real} rate is based on constant purchasing power: $1 + i_R = (1 + i_A)/(1 + f)$
    \item Actual values use the actual MARR; real values use the real MARR
    \item Always use the \b{actual} MARR in tax benefit factor calculations
\end{enumerate}
